% Abstract — plain English first; define-before-use; headline numbers in one paragraph.

A pension fund manager has two jobs at once: do not lose more than a set percentage from the portfolio's recent peak, and still earn more than cash. That first job is called controlling \emph{drawdown}. Standard tools make you pick one job or the other. Mean-variance optimisation~\parencite{markowitz1952portfolio} picks weights once and ignores the running peak. Trailing stop-loss rules sell only \emph{after} a fall has already happened. Value-at-Risk~\parencite{rockafellar2000optimization} measures one bad day at a time and says nothing about a long slide from peak to trough.

This dissertation asks whether a small AI agent can do better on both jobs. The agent learns day-by-day buy and sell decisions through \emph{reinforcement learning} --- trial and error in a simulated market. The specific algorithm is Proximal Policy Optimization (PPO)~\parencite{schulman2017ppo}, a standard method for sequential decisions; it is used unchanged here. The agent is paired with a \emph{forecaster}: a neural network that reads recent price history and outputs tomorrow's expected return \emph{and} how wide its prediction interval is. That interval width becomes a \textbf{confidence signal} (formally the uncertainty score $u_t$): near zero means the model is confident; near one means it is not. When the signal is high, the agent shrinks trades or blocks new buys.

The agent is tested against three named alternatives on the same fixed protocol: passive buy-and-hold, a manual trailing stop-loss rule, and the same PPO agent \emph{without} the confidence signal. Training and test dates, random seeds, and performance measures are identical across all four so the comparison is fair.

The headline result uses 70 US stocks and ten independent training runs per stock. On the held-out 2022--2025 window, the confidence-aware agent grows a \$1\,million start to a \emph{median} of about \$1.6\,million. Its median Sharpe ratio (return per unit of risk; higher is better) is $+0.68$ to $+0.70$, and it finishes ahead of its starting capital on 89--92\,\% of stocks. The plain PPO control finishes flat at about \$999\,000, with a median Sharpe of $-0.04$ and a win-rate of 46\,\%. A standard paired statistical test on the 70 matched stocks rejects ``no difference'' at $p < 10^{-12}$; a bootstrap confidence interval puts the median gain at \$559k to \$814k.

The pattern holds on earlier data the agent never trained on. A four-fold walk-forward check slides the test window through 2018--2025, covering the COVID crash and the 2022 inflation shock. The confidence-aware agent wins every fold (87\,\% win-rate overall; median Sharpe $+0.55$; $p < 10^{-8}$).

A follow-up arm swaps \emph{how} the confidence signal is computed: one version reads market noise (aleatoric uncertainty), another reads model unfamiliarity via Monte-Carlo dropout~\parencite{gal2016dropout} (epistemic uncertainty). The two versions land close on final wealth (difference not significant at the 5\,\% level) and differ only slightly on Sharpe ($+0.014$ in favour of the noise-based version, $p = 7 \times 10^{-4}$). Any well-calibrated confidence signal helps; the exact source is a second-order choice.

The contribution is a careful empirical comparison, not a new algorithm. What is new is the four-way test on one reproducible protocol --- including the trailing-stop comparator that much of the finance reinforcement-learning literature omits --- and a public notebook that reproduces the pipeline end to end.

\vspace{1em}
\noindent\textbf{Keywords:} portfolio protection; confidence signals; drawdown control; reinforcement learning; walk-forward validation.
